\documentclass{article}
\usepackage[UTF8]{ctex}
\usepackage{amsmath,amssymb,amsfonts}
\usepackage{enumitem}
\usepackage{geometry}
\geometry{a4paper, margin=1in}
% === 错题本排版优化（不参与题目标识） ===
\setlength{\parindent}{0pt}
\setlength{\parskip}{0.35em plus 0.10em minus 0.08em}
\setlength{\emergencystretch}{3em}
\setlist{itemsep=0.25em, topsep=0.40em, parsep=0pt}
\allowdisplaybreaks[2]
% === 错题本排版优化结束 ===
\begin{document}
\section*{上海财经大学2025年高等代数真题}
\section*{44．上海财经大学2025年高等代数试卷（回忆版）}

\subsection{题目1}
1．已知 \(A\) 为 \(n\) 阶矩阵，证明：\(A^{3}=I\) 的充要条件是 \(r(A-I)+r\left(A^{2}+A+I\right)=n\) ．

\subsection{题目2}
2．求行列式

\[
D_{n}=\left|\begin{array}{ccccc}
0 & 1 & 1 & \cdots & 1 \\
1 & 0 & x & \cdots & x \\
1 & x & 0 & \cdots & x \\
\vdots & \vdots & \vdots & & \vdots \\
1 & x & x & \cdots & 0
\end{array}\right|
\]

\subsection{题目3}
3．设 \(A, B\) 都是 \(n\) 阶复矩阵，\(f_{A}(x), f_{B}(x)\) 分别是它们的多项式。证明下列说法等价：

\begin{itemize}
  \item[（1）] \(A, B\) 没有公共的特征值．
  \item[（2）] \(\left(f_{A}(x), f_{B}(x)\right)=1\) ．
  \item[（3）] \(A X=X B\) 只有零解．
\end{itemize}

\subsection{题目4}
4．\(V\) 是复数域上所有 \(n\) 阶矩阵的全体组成的线性空间， \(\mathcal{A}\) 是 \(V\) 上的线性变换， \(\mathcal{A}(X)=A X-X A^{\mathrm{T}}\)（注：另外一个版本是 \(\mathcal{A}(X)=A X\) ），证明：如果 \(A\) 可对角化，则 \(\mathcal{A}\) 可对角化。

\subsection{题目5}
5．实二次型 \(f\left(x_{1}, x_{2}, x_{3}\right)=x_{1}^{2}+x_{2}^{2}+x_{3}^{2}+2 a x_{1} x_{2}+2 x_{1} x_{3}+2 b x_{2} x_{3}\) 经过正交变换 \(X=P Y\) 得标准形 \(f=y_{2}^{2}+2 y_{3}^{2}\) ，求 \(a, b\) 的值和正交矩阵 \(P\) 。

\subsection{题目6}
6．已知

\[
A=\left(\begin{array}{llll}
1 & 2 & 0 & 0 \\
2 & 5 & 0 & 0 \\
0 & 0 & 2 & 2 \\
4 & 0 & 2 & 1
\end{array}\right)
\]

矩阵 \(X\) 满足 \(\left(A^{*}\right)^{*} X A^{*}=8(A X-A)\) ，求 \(A^{*},\left(A^{*}\right)^{*},\left|\left(A^{*}\right)^{*}\right|, X\) 。

\subsection{题目7}
7．\(A, B, C\) 均为实矩阵，证明：\(\left(\begin{array}{cc}A & B \\ B^{\mathrm{T}} & C\end{array}\right)\) 正定的充要条件是 \(A\) 与 \(C-B^{\mathrm{T}} A^{-1} B\) 正定。

\subsection{题目8}
8．给定方程组

\[
\left\{\begin{array}{l}
x_{1}-x_{2}-x_{3}+x_{4}=1 \\
x_{1}-3 x_{2}-3 x_{3}+x_{4}=1
\end{array}\right.
\]

\begin{itemize}
  \item[（1）] 求方程组的通解。
  \item[（2）] 最多可以找出多少个线性无关的解向量？找出一组用它来表示所有的解．
\end{itemize}

\subsection{题目9}
9．求矩阵 \(A=\left(\begin{array}{ccc}3 & 1 & 0 \\ -1 & 1 & 0 \\ -1 & -1 & 2\end{array}\right)\) 的 Jordan 标准形，并计算 \(A^{n}+A^{n-1}+I\) ．

\subsection{题目10}
10．设 \(V\) 是复数域上的 \(n\) 维线性空间，\(\varphi, \psi\) 为 \(V\) 上的线性变换，且 \(\varphi^{2}=\psi^{2}=\operatorname{id}_{V}\) ，这里 \(\operatorname{id}_{V}\) 是 \(V\)上的恒等变换。证明：

\begin{itemize}
  \item[（1）] 若 \(n\) 是奇数，则 \(\varphi\) 和 \(\psi\) 有公共的特征向量。
  \item[（2）] 若 \(n\) 是偶数，则存在一子空间 \(W\) 同时是 \(\varphi\) 和 \(\psi\) 的不变子空间，且 \(\operatorname{dim} W=1\) 或 2 。
\end{itemize}
\end{document}
